\documentclass[../marker.tex]{subfiles}
\begin{document}
\part{LỘ TRÌNH}
\begin{tomtat}
Phần E chỉ có một chương và nó không dạy khái niệm mới. Nó trả lời câu hỏi thực dụng nhất: \emph{nếu tôi phải làm xong hệ thống docking này, tôi làm gì trước, làm gì sau, và vì sao đúng thứ tự ấy?} Lộ trình gồm tám bước, và điểm đáng chú ý của nó là ba bước đầu không có dòng code nào --- chúng là tính toán trên giấy và quyết định cơ khí. Nếu chỉ nhớ một điều: \emph{bước đầu tiên là viết script Monte Carlo cho constellation dự kiến, trước khi mua camera và trước khi đặt gia công tấm marker}. Một buổi chiều ở bước ấy quyết định phần lớn kết quả cuối cùng.
\end{tomtat}
\subfile{00-map}
\subfile{19-lo-trinh-toi-thieu/lo-trinh-toi-thieu}
\ifSubfilesClassLoaded{\printbibliography[title={Nguồn}]}{}
\end{document}
